\documentclass[11pt]{article}

\pagestyle{myheadings}

\usepackage{times}
\usepackage{color}
\usepackage{alltt}
\usepackage{graphicx}

%% Declarations of formatting parameters & terminology (HANDOUT VERSION)
\raggedbottom

%% Page layout
\oddsidemargin 0pt
\evensidemargin 0pt

\textheight 600pt
\textwidth 469pt

\setlength{\parskip}{1ex}
\setlength{\parindent}{0em}

%%% Tag to indicate references outside of handout %%%%
%%% Argument should be something like Chapter or Section 
\newcommand{\csappiii}[1]{CS:APP3e #1}

%%%%%%%%%%%%%%%%%%%%%%%%%%%%%%%%%%%
% Environments for handout problems
%%%%%%%%%%%%%%%%%%%%%%%%%%%%%%%%%%%

%% For problems.  Arguments are a name and a category rating
\newcounter{prob}
\newenvironment{problem}[2]{%
\refstepcounter{prob}%
\begin{center}%
\begin{minipage}{5.75in}%
\label{#1}%
{\bf Problem \arabic{prob}}
[Category #2]:\par}{\medskip\end{minipage}\end{center}}


%% For practice problems.  Arguments are a name and a category rating
\newenvironment{pproblem}[2]{%
\refstepcounter{prob}%
\begin{quote}%
\label{#1}%
{\bf Practice Problem \arabic{prob}}:\par\small}{\medskip\end{quote}}

%% For homework problems.  Arguments are a name and a category rating
\newenvironment{hproblem}[2]{%
\refstepcounter{prob}%
\label{#1}%
{\bf Homework Problem \arabic{prob}}
\showcat{#2}:\par}{\mbox{}\par}

%% For solutions.  Argument is label given to problem
\newenvironment{solution}[1]{%
{\bf Problem \ref{#1} Solution: [Pg.~\pageref{#1}]}\par}{\medskip}

%% Generate Alphabetically-numbered subproblems
\newenvironment{subproblem}{\begin{enumerate}%
\def\labelenumi{\Alph{enumi}.}}{\end{enumerate}}



%% Generic environments and commands that we use in the CS:APP book

%%%%%%%%%% 
% Commands
%%%%%%%%%% 

% Reference to versions of book
\newcommand{\csapp}{CS:APP}
\newcommand{\csappii}{CS:APP2e}

% Reference to something in book from a web aside
\newcommand{\bookref}[1]{CS:APP3e-\ref{#1}}
% Reference to a named web aside
\newcommand{\wasideref}[1]{\textsc{#1}}

% For drawing different kinds of circles in the course table in the preface
% fcirc: filled circle, % ecirc: empty circle
\newcommand{\fcirc}{\begin{picture}(1,1)\put(1,1){\circle*{5}}\end{picture}}
\newcommand{\ecirc}{\begin{picture}(1,1)\put(1,1){\circle{5}}\put(1,1){\circle*{2.4}}\end{picture}}

% For putting labels on pictures
\newcommand{\clabel}[1]{\makebox(0,0){#1}}
\newcommand{\rlabel}[1]{\makebox(0,0)[r]{#1}}
\newcommand{\llabel}[1]{\makebox(0,0)[l]{#1}}
\newcommand{\blabel}[1]{\makebox(0,0)[b]{#1}}
\newcommand{\lblabel}[1]{\makebox(0,0)[lb]{#1}}
\newcommand{\tlabel}[1]{\makebox(0,0)[t]{#1}}

%% Comment command
\newcommand{\comment}[1]{}

%% Two part captions.  First part is title.  Second is explanation.
\newcommand{\mycaption}[2]{\caption[#1]{\small \textsf{\textbf{#1} #2}}}

%% For showing categories of homework problems
\definecolor{catcolor}{rgb}{0,0,1}
\newcommand{\catmark}{{\color{catcolor}\blacklozenge}}
\newcommand{\catone}{$\catmark$}
\newcommand{\cattwo}{$\catmark\,\catmark$}
\newcommand{\catthree}{$\catmark\,\catmark\,\catmark$}
\newcommand{\catfour}{$\catmark\,\catmark\,\catmark\,\catmark$}
\newcommand{\showcat}[1]{\ifthenelse{#1=1}{\catone}{%
 \ifthenelse{#1=2}{\cattwo}{%
   \ifthenelse{#1=3}{\catthree}{\catfour}}}}


% Linker ref->def notation
\newcommand{\refdef}[2]{$\texttt{REF(#1)} \rightarrow \texttt{DEF(#2)}$}

% Linker ref->def question
\newcommand{\refdefq}[1]{$\texttt{REF(#1)} \rightarrow \texttt{DEF(\_\_\_\_\_\_\_\_\_\_.\_\_\_\_)}$}

%%%%%%%%%%%%%%%%%%%%%%%%%%%%%%%%%%%%%%%%%%%%%%%%%%%%%%%%%%%% 
% Environments and macros for different elements of the book
%%%%%%%%%%%%%%%%%%%%%%%%%%%%%%%%%%%%%%%%%%%%%%%%%%%%%%%%%%%%% 

%% This TeX counter is used to generate references to 
%% different code line numbers
\newcounter{codelineno}

%% This defines the color we are using for all the code line numbers
\definecolor{linenumcolor}{rgb}{0,0,1} 

%% Various Unix command-line prompts
\newcommand{\unixprompt}{linux>}
\newcommand{\linuxprompt}{linux>}
\newcommand{\solarisprompt}{solaris>}

%% Definition of RIO, SIO  package names
\newcommand{\rio}{\textsc{Rio}}
\newcommand{\sio}{\textsc{Sio}}
\newcommand{\Tiny}{\textsc{Tiny}}

%%%
%%% Environments
%%%

%% pcode- Generating pseudo-code
\newenvironment{pcode}{\mbox{}\newline\vspace{.5ex}\newline\begin{minipage}{6in}\begin{tt}\begin{tabbing}MMM\=MMM\=MMM\=MMM\=MMM\=MMM\=MMM\=MMM\=MMM\=\+\kill}{\end{tabbing}\end{tt}\end{minipage}\newline\vspace{.5ex}\newline}

%% Showing syntactic elements in pseudo-code
\newcommand{\syntax}[1]{{\rm\it #1}}

%% aside - for displaying sidebars 
\newenvironment{aside}[1]%
{\begin{quote}\footnotesize{\textbf{Aside: #1}}\\}
{\textbf{End Aside.}\end{quote}}

%% waside - for highlighting content available on the WEE
%% Arg1: Code.  Arg2: Title
\newenvironment{waside}[2]%
{\begin{quote}\footnotesize{\textbf{Web Aside \textsc{#1}: #2}}\\}
{\textbf{End Aside.}\end{quote}}

%% ntc - New to C element
\newenvironment{ntc}[1]%
{\begin{quote}\footnotesize{\textbf{New to C?: #1}}\\}
{\textbf{End.}\end{quote}}

%% principle - to highlight key concept
\newenvironment{principle}[1]%
{\begin{quote}\textbf{Principle}: {\em #1} \hrulefill \mbox{} \vskip 1pt}%
{\vskip 1pt \hrulefill \mbox{} \end{quote}}

%% derivation - to present derivation
\newenvironment{derivation}[1]%
{\begin{quote}\textbf{Derivation}: {\em #1} \hrulefill \mbox{} \vskip 1pt}%
{\vskip 1pt \hrulefill \mbox{} \end{quote}}

%% ccode- for displaying formatted C code (c2tex) 
\newenvironment{ccode}%
{\small}%
{}

%% scode - for displaying formatted ASM code (s2tex and d2tex)
\newenvironment{scode}%
{\small}
{}

%% codefrag - for displaying unformatted code fragments
\newenvironment{codefrag}%
{\small\begin{alltt}}%
{\end{alltt}%
}

%% tty - for displaying TTY input and output
\newenvironment{tty}%
{\small\begin{alltt}}%
{\end{alltt}}


%% Environment for descriptions that provide aligned text
\newenvironment{mydesc}[1]{
\setbox1=\hbox{#1}
\begin{list}{}{
\setlength{\labelwidth}{\wd1}
\setlength{\leftmargin}{\wd1}
  \addtolength{\leftmargin}{1em}
  \addtolength{\leftmargin}{\labelsep}
\setlength{\rightmargin}{1em}}}{\end{list}}

\newcommand{\litem}[1]{\item[#1\hfill]}
\newcommand{\ritem}[1]{\item[#1]}

%% For marking different verions of problems

\newcommand{\sameprob}[1]{\textbf{Same as Version A Problem \ref{#1}}\\ \mbox{}}

%%%%%%%%%%%%%%%%%%%%%%%%%%%%%%%%%%%%%%%%%%%%%%%%%%%%%%%%%%%%%%%%%%%%%%%%%%%%
%% Indexing 
%%%%%%%%%%%%%%%%%%%%%%%%%%%%%%%%%%%%%%%%%%%%%%%%%%%%%%%%%%%%%%%%%%%%%%%%%%%%
%% Distinguish multiple types of index entries
%%   - IA32 instructions (iindex, diindex)
%%   - x86-64 instructions (xindex, dxindex)
%%   - Y86  instructions (yindex, dyindex)
%%   - HCL  things (hindex, dhindex)
%%   - Program names (pindex, dpindex)
%%   - C things (cindex, dcindex)
%%   - Stdlib functions (sindex, dsindex)
%%   - Unix functions (uindex, duindex)
%%   - CSAPP functions (csindex, dcsindex)
%%   - CSAPP programs (pcsindex, dpcsindex)
%%   - Library files, such as <limits.h> (lindex, dlindex) 
%%   - Everything else. (index, dindex)
%% In every case, the `d' version is used for the defining index entry

%% Generic instructions.
%% Takes two arguments:
%% The name of the instruction
%% An English rendition of the mnemonic
\newcommand{\gindex}[2]{\index{#1@\textsc{#1} [Instruction class] #2}}
\newcommand{\dgindex}[2]{\index{#1@\textsc{#1} [Instruction class] #2|emph}}


%% IA32 instructions.
%% Takes two arguments:
%% The name of the instruction
%% An English rendition of the mnemonic
\newcommand{\iindex}[2]{\index{#1@\texttt{#1} [IA32] #2}}
\newcommand{\diindex}[2]{\index{#1@\texttt{#1} [IA32] #2|emph}}

%% x86-64 instructions.
%% Takes two arguments:
%% The name of the instruction
%% An English rendition of the mnemonic
\newcommand{\xindex}[2]{\index{#1@\texttt{#1} [x86-64] #2}}
\newcommand{\dxindex}[2]{\index{#1@\texttt{#1} [x86-64] #2|emph}}

%% y86-64 instructions.
%% Takes two arguments:
%% The name of the instruction
%% An English rendition of the mnemonic
\newcommand{\yindex}[2]{\index{#1@\texttt{#1} [Y86-64] #2}}
\newcommand{\dyindex}[2]{\index{#1@\texttt{#1} [Y86-64] #2|emph}}

%% HCL things
%% Takes two arguments:
%% The name of the instruction
%% An English rendition of the mnemonic
\newcommand{\hindex}[2]{\index{#1@\texttt{#1} [HCL] #2}}
\newcommand{\dhindex}[2]{\index{#1@\texttt{#1} [HCL] #2|emph}}

%% Program names (e.g., gcc, gdb.  Write in lower case)
%% Takes two arguments:
%% The name of the program
%% An English description
\newcommand{\pindex}[2]{\index{#1@{\sc #1} #2}}
\newcommand{\dpindex}[2]{\index{#1@{\sc #1} #2|emph}}

%% C things (operators, statements, etc)
%% Takes two arguments:
%% C name
%% An English description
\newcommand{\cindex}[2]{\index{#1@\texttt{#1} [C] #2}}
\newcommand{\dcindex}[2]{\index{#1@\texttt{#1} [C] #2|emph}}

%% Stdlib functions
%% Takes two arguments:
%% function name
%% An English description
\newcommand{\sindex}[2]{\index{#1@\texttt{#1} [C Stdlib] #2}}
\newcommand{\dsindex}[2]{\index{#1@\texttt{#1} [C Stdlib] #2|emph}}

%% Unix functions
%% Takes two arguments:
%% function name
%% An English description
\newcommand{\uindex}[2]{\index{#1@\texttt{#1} [Unix] #2}}
\newcommand{\duindex}[2]{\index{#1@\texttt{#1} [Unix] #2|emph}}

%% CSAPP functions
%% Takes two arguments:
%% function name
%% An English description
\newcommand{\csindex}[2]{\index{#1@\texttt{#1} [CS:APP] #2}}
\newcommand{\dcsindex}[2]{\index{#1@\texttt{#1} [CS:APP] #2|emph}}

%% CSAPP programs
%% Takes two arguments:
%% function name
%% An English description
\newcommand{\pcsindex}[2]{\index{#1@{\sc #1} [CS:APP] #2}}
\newcommand{\dpcsindex}[2]{\index{#1@{\sc #1} [CS:APP] #2|emph}}

%% Library file names (e.g., <limits.h>)
%% Takes two arguments:
%% The name of the file (without brackets)
%% An English description
\newcommand{\lindex}[2]{\index{#1@\texttt{<#1>} #2}}
\newcommand{\dlindex}[2]{\index{#1@\texttt{<#1>} #2|emph}}

%% Normal index entry.
%% Takes one argument
%% \index already defined
\newcommand{\dindex}[1]{\index{#1|emph}}





\newcommand{\pxydrive}{\textsc{PxyDrive}}
\newcommand{\pxyregress}{\textsc{PxyRegress}}

\begin{document}


\title{
\includegraphics[width=4in]{figs/pixie-driver} \\[2ex]
\pxydrive: A Proxy Testing Framework\thanks{Copyright \copyright{} 2018, R. E. Bryant.  All rights reserved.}
}

\author{Randal E. Bryant \\ Carnegie Mellon University}

\maketitle

\section{Overview}
\label{sec:overview}

\begin{figure}
\begin{center}
\includegraphics[scale=0.6]{figs/structure}
\end{center}
\mycaption{Structure of \pxydrive}{The program acts as both the client and the servers for a proxy program.}
\label{fig:structure}
\end{figure}

\pxydrive{} (pronounced ``pixie drive'') provides a comprehensive
testing environment for web proxy programs.  It enables testing and
debugging of the basic operations of a proxy, as well as the ability
to support concurrent requests and to provide caching of retrieved
files.  An associated program \pxyregress{} runs a series of standard
tests, using \pxydrive{} to perform each test.  These two programs are
designed to help students implement their proxies and to help graders
perform automated testing.  \pxydrive{} was developed to support the
ProxyLab assignment, but it can be used on other proxies, as well.

\newpage

Figure \ref{fig:structure} illustrates the overall structure of
\pxydrive{}.  As shown, it includes a number of components:
\begin{itemize}
\item An external executable {\em proxy} program.  It is executed as child
  process of \pxydrive.
\item A web {\em client} that can issue requests for files and then compare
  the responses to original copies of the files.
\item One or more {\em servers} that can respond to requests for files
\item A set of {\em files} shared by the client and the servers
\item A {\em console} providing a simple command-line interpreter to
  perform tests and to observe results.
\end{itemize}

This structure provides a number useful features for testing a proxy
program:
\begin{itemize}
\item By including all of these components in a single program, \pxydrive{}
has very detailed control over the sequencing of interactions with the
proxy.
\item By having files shared between the client and the servers, the
  client can check that the results it receives match those sent by
  the servers.
\item By observing a transaction from both the client's and the
  server's side, it can provide useful information to the user about
  all of the communications that take place.
\item The ability to observe both server and client communications
  enables \pxydrive{} to check that the proxy complies with any
  formatting requirements for HTTP messages.
\end{itemize}

\section{Running \pxydrive}
\label{sec:running}

\pxydrive{} is written in Python, and so is invoked as
\texttt{./pxydrive.py} from its local directory.  (You can also run it
from another directory by giving an appropriate path.)  \pxydrive{}
has been tested on both Linux and Mac~OS~X.  Its ability to run on
other platforms is unknown.

\subsection*{Command-Line Arguments}

\pxydrive{} supports the following command-line
arguments:
\begin{description}
\item[\texttt{-h}]  $\;$ \\
Print usage information
\item[\texttt{-p} {\it PROXY}] $\;$ \\
Run specified proxy program.  For
  more details, consult the documentation of the \texttt{proxy}
  command.
\item[\texttt{-f} {\it FILE}]  $\;$ \\
Read commands from {\it FILE}\@.
For more details, consult the documentation of the
  \texttt{source} command.
\item[\texttt{-l} {\it FILE}]  $\;$ \\
Write inputs and outputs to {\it FILE}\@.
For more details, consult the documentation of the
  \texttt{log} command.
\item[\texttt{-S}] $\;$ \\
Enforce strict formatting rules for HTTP request message headers. 
For more details, consult the documentation of the \texttt{strict} option.
\end{description}

\pxydrive{} creates two subdirectories of the directory in which it is
invoked:
\verb@source_files@, in which generated files are stored, and \verb@response_files@,
  in which files received by the client are stored.  These files can
  be examined when some transaction generates an error.
  Old versions of
  these files are deleted every time \pxydrive{} starts.

Once started, the user directs \pxydrive{} to perform a series of
tests, either by typing commands directly to the program or by reading
the commands from a file.

\subsection*{Proxy Requirements}

\pxydrive{} is designed to tests proxy programs that meet the
requirements specified in ProxyLab.  The lab description specifies several important aspects of the request headers sent by the proxy to a server:
\begin{itemize}
\item There must be a \texttt{Host} header
\item There must be \texttt{User-Agent}, \texttt{Connection} and \texttt{Proxy-Connection} headers with specific values.
\item Any other request headers that are sent by the client should be forwarded unchanged.
\end{itemize}
The latter requirement is especially important, because \pxydrive{} uses information passed in headers
 to track the interactions between clients and servers.

\pxydrive{} can operate in two different modes.  In {\em strict} mode
(set with either the command-line argument \texttt{-S} or the \texttt{strict} option), these
requirements are strictly enforced.  Error messages will be printed for request messages that do not comply.
In {\em non-strict} mode (the
default), the program doesn't require these headers.  Instead, it does
its best to track the interactions heuristically.  Even in non-strict
mode, it works best when the proxy
forwards information placed in the request headers to the server.

\section{\pxydrive{} Operation}
\label{sec:operation}


\begin{figure}
\begin{center}
\includegraphics[scale=0.6]{figs/transaction}
\end{center}
\mycaption{Steps in a transaction}{A transaction involves a {\em request} (steps 1--3) and a {\em response} (steps 4--6)}
\label{fig:transaction}
\end{figure}


A {\em transaction} by \pxydrive{} consists of all of the steps required for the
client to get a file from a server via the proxy and ascertain its correctness.  A transaction
has two different phases: the {\em request} by the client,
followed by the {\em response} by the server.  Each transaction can be
divided into six steps---three for each phase, as is illustrated in
Figure \ref{fig:transaction}:
\begin{description}
\item[Request] $\;$ \\
\begin{enumerate}
\item The client sends a \texttt{GET} request to the proxy.  This
  request includes a URL, identifying the host server and the file name.
\item The proxy forwards the request to the designated server.
\item The server receives the request, locates the file, and
  formulates a response message.
\end{enumerate}
\item[Response] $\;$ \\
\begin{enumerate}
\setcounter{enumi}{3}
\item The server sends the response message to the proxy.
\item The proxy forwards the response to the client.
\item The client receives the requested file and compares it to the
  original one.
\end{enumerate}
\end{description}

\begin{figure}
\begin{center}
\input{examples/A01}
\end{center}
\mycaption{Using \pxydrive{} to perform transaction with separate request and response.}{Delays must be included to allow the proxy, the server, and the network to operate.}
\label{cmd:A01}
\end{figure}

\pxydrive{} gives fine-grained control over transactions, providing
separate commands to trigger the request phase and the response
phase.  This feature makes it possible to test important features of
the proxy, including its ability to handle concurrent transactions and
to cache previously retrieved files.

The following simple examples illustrate how the user can control
\pxydrive{} operation.  These examples all consist of commands that
can be supplied to the command-line interface.  Observe that
`\verb@#@' is a comment character---it and anything to the right on
the same line are ignored.

Figure \ref{cmd:A01} shows a session having a single transaction, with
separate commands to trigger the request and the response.  (For this
session, the proxy was specified as a command-line option.)
The lines are as follows:
\begin{enumerate}
\item A server is set up and given name \texttt{s1}.  (It is possible
  to have multiple servers, each with a unique name.)
\item A file, consisting of 10,000 random ASCII characters is
  generated.  (The extension \texttt{.txt} indicates that the file
  should contain only ASCII characters.)  
The suffix `\texttt{K}' indicates that the
  number of bytes should be multplied by 1,000.
\item The client makes a request for this file from the server.  The transaction is
  given the identifier `\texttt{r1}.'  Transaction identifiers must be
  unique, because they provide the mechanism by which \pxydrive{}
  tracks the progress of transactions.
\item The program waits for 100 milliseconds, giving time for steps
  1--3 of the transaction to occur.  Explicit delays must be included
  whenever the client, proxy, or server must complete some operations
  before the next command can be performed.
\item The server is instructed to respond to the request.  Note the
  use of the transaction identifier to indicate the specific request.
\item The program waits for 100 milliseconds, giving time for steps
  4--6 of the transaction to occur.
\item
The program checks that the file was received correctly by the
client.  An error message will be printed if either the server could
not find the file, the file was corrupted, or the client did not
receive the reponse.
\end{enumerate}

\begin{figure}
\begin{center}
\input{examples/A02}
\end{center}
\mycaption{Using \pxydrive{} to perform transaction with immediate response.}{Only one delay is required.}
\label{cmd:A02}
\end{figure}

Figure \ref{cmd:A02} also performs a single transaction, but it
illustrates several additional features of \pxydrive{}.  In line 2,
the generated file is given the extension \texttt{.bin}.  This causes
the generated file to contain arbitrary bytes.  Such a file will test
the ability of the proxy to properly handle binary data.  In addition
the \texttt{fetch} command of line 3 indicates that \pxydrive{} should
perform a complete transaction, with the server responding immediately
to the request.  Only one delay command (line 4) is required, since all six
steps of the transaction will take place.
Running a session consisting of a mix of fetches and
requests requires the proxy to handle a variety of event timings.


\begin{figure}
\begin{center}
\input{examples/A03}
\end{center}
\mycaption{Using \pxydrive{} to perform transaction with missing file.}{The check confirms status code 404.}
\label{cmd:A03}
\end{figure}

Figure \ref{cmd:A03} illustrates how \pxydrive{} can be used to test
the ability of the proxy to handle transactions involving exceptional
conditions.  It shows an attempt to retrieve a file that does not
exist.  The server will respond to this request with a message having
status code 404, indicating that the file was not found.  The proxy
should forward this message to the client.  The optional argument
\texttt{404} on line 4 indicates that the program should check for
this status code in the response to the client.

These three examples illustrate some of the capabilities of
\pxydrive{} to run tests of a proxy.  Section \ref{sec:testing} provides additional examples
showing how to test different features of a proxy.

\section{Commands}
\label{sec:command}

The following documents all of the commands supported by \pxydrive{}.
These are divided into different categories, according to when and how
they get used.

In the following, square brackets enclose optional arguments, while
the notation \textit{ARG}$+$ means one or more repetitions of argument \textit{ARG}\@.
The notation (\textit{ARGA}$|$\textit{ARGB}) indicates a choice
between arguments \textit{ARGA} and \textit{ARGB}\@.

\subsection*{Set Up}

The following commands are used at the beginning of a session to
establish the environment for a set of transactions.

\begin{description}
\item[\texttt{proxy} {\it PATH} [\textit{ARG}$+$]] $\;$ \\
    Execute proxy program stored as executable file {\it PATH}\@.
    Any remaining arguments are passed to the proxy as command-line
    arguments.  (There is no need to pass a port number---this is
    supplied by \pxydrive{}.)
    Any results printed by the proxy to either standard
    output or standard error will be echoed to the \pxydrive{}
    standard output
    and log files.  If some proxy is already running, it will be
    terminated before starting the specified proxy.

\item[\texttt{serve} \textit{SID}$+$] $\;$ \\
    Set up server(s) with named identifier(s).  These will be assigned
    ports chosen by \pxydrive{}.  Server identifiers must be unique.

\item[\texttt{generate} \textit{NAME}\texttt{.}(\texttt{txt}$|$\texttt{bin}) \textit{BYTES}] $\;$ \\
    Generate file with random content.  
    Files with names of the form \textit{NAME}\texttt{.txt} will
    contain random ASCII characters, formatted into lines of length
    80.
    Files with names of the form \textit{NAME}\texttt{.bin} will
    contain random, unformatted bytes.
    Argument \textit{BYTES} indicates the number of bytes in the
    file.  As a shortcut, suffix `\texttt{K}' multiplies the number of
    bytes by 1,000, and suffix `\texttt{M}' multiplies the number of
    bytes by 1,000,000.
\end{description}

\subsection*{Transactions}

The following commands control the progress of transactions and
observe their results.

\begin{description}
\item[\texttt{fetch} \textit{ID} \textit{FILE} \textit{SID}]  $\;$ \\
    Fetch file from server (i.e., the request will be passed to
    server, and the server will be allowed to respond immediately.)
    \textit{ID} is an identifying name for the transaction.  It must
    be unique.  Unless testing for a ``not found'' response, \textit{FILE} should be one of the files accessible to
    the server (e.g., one generated with the \texttt{generate} command).
    \textit{SID} identifies which server to use.

\item[\texttt{request} \textit{ID} \textit{FILE} \textit{SID}]  $\;$ \\
    Request file from server (i.e., the request will be passed to
    server, but the server will be not be allowed to respond.)
    \textit{ID} is an identifying name for the transaction.  It must
    be unique.  Unless testing for a ``not found'' response, \textit{FILE} should be one of the files accessible to
    the server (e.g., one generated with the \texttt{generate} command).
    \textit{SID} identifies which server to use.

\item[\texttt{respond} \textit{ID}$+$]  $\;$ \\
    Allow the server(s) to respond to request(s) specified by the
    identifier(s).  An error message will be printed if the server has
    not yet received this request.

\item[\texttt{check} \textit{ID} [\textit{CODE}]]  $\;$ \\
    Check that transaction \textit{ID} generated the proper status
    code.  The default code is 200, indicating that the file was
    successfully retrieved.    Additional status
codes include 400 (bad request), 404 (not
found), 501 (not implemented), and 505 (unsupported HTTP version), as
are documented in Figure 11.25 of CS:APP\@.

\item[\texttt{trace} \textit{ID}$+$]  $\;$ \\
    Trace the histories of the specified transactions.
    This will print information about the six transaction steps illustrated in
    Figure \ref{fig:transaction}, including copies of the HTTP
    headers. This command can be very useful when debugging a proxy or
    a test sequence.
\end{description}

\subsection*{Runtime Control}

The following commands control the operation of the program.

\begin{description}
\item[\texttt{delay} \textit{MS}]  $\;$ \\
    Delay for \textit{MS} milliseconds.

\item[\texttt{log} \textit{FILE}]  $\;$ \\
    Copy outputs to \textit{FILE}\@.  Inputs are also echoed to the file
    if the \texttt{echo} option is set. (This is the default.)

\item[\texttt{source} \textit{FILE}]  $\;$ \\
    Read commands from \textit{FILE}\@.

\item[\texttt{option} \textit{OPTION} \textit{VALUE}]  $\;$ \\
    Set a runtime option.  Supported options are documented below.

\item[\texttt{delete} \textit{FILE}]  $\;$ \\
    Delete file \textit{FILE}\@.  Files generated during previous runs
    of \pxydrive{} are automatically deleted every time the program
    starts, and so explicit deletion of files is only required for
    running special tests, or to avoid leaving large files around
    after a test has completed.

\item[\texttt{help}] $\;$ \\
Print a summary of the commands and options.

\item[\texttt{quit}]  $\;$ \\
Exit the program.  This may take awhile, since multiple threads must
be terminated.  If the program hangs up indefinitely or causes an
unexpected exception, it may be necessary to explicitly kill the
\pxydrive{} process.
\end{description}

\subsection*{Options}

The following options are supported.  For Boolean options, 1
indicates true, and 0 indicates false.
\begin{description}
\item[autotrace] (Boolean, default = 0)\\
 Trace every request for which a check fails.
This mode is helpful for debugging,
because it causes detailed information about each failed check to be
printed.
\item[echo] (Boolean, default = 1)\\
Echo inputs to log files.  This makes it possible to keep track of which commands were being executed.
\item[error] (Integer, default = 5)\\
The maximum number of errors that can occur before the program aborts.
\item[stretch] (Integer, default = 100):\\
Adjusts the length of each
  delay.  Given a command ``\textbf{delay} $d$,'' and given a stretch
  value of $s$, the actual delay will be $d \cdot s/100$ milliseconds.
  Setting the stretch to a value greater than 100 lengthens all
  delays, while setting it to a value less than 100 shortens them.
  This provides a way to adjust all of the delays in a test sequence
  without editting all of the delay commands in a file.  It can be
  useful in detecting cases where tests are just on a border between
  success and failure due to event timings.
\item[strict] (Boolean, default = 0):\\ Apply strict formatting rules
  to HTTP request headers.  As described in the ProxyLab writeup,
  request messages are required to include headers with settings for
  \texttt{Host}, \texttt{Connection}, and \texttt{Proxy-Connection}.
  In addition, the proxy should pass along any other headers defined
  in the request message from the client.
\end{description}

\section{Testing Proxy Capabilities}
\label{sec:testing}

The following are some examples demonstrating how the fine-grained
control provided by \pxydrive{} allows testing of different features
of a proxy.  Additional examples can be found in the standard tests
described in Section \ref{sec:standard}.

\begin{figure}
\begin{center}
\input{examples/B01}
\end{center}
\mycaption{Using \pxydrive{} to perform transaction requiring concurrent proxy.}{A serial proxy would be held up waiting for the server to respond to request \texttt{r1}.}
\label{cmd:B01}
\end{figure}

Figure \ref{cmd:B01} illustrates the ability of \pxydrive{} to
ascertain that a proxy supports concurrent transactions.  In this
test, two requests---\texttt{r1} and \texttt{r2}---are made (lines 4--5), but the
server only responds to \texttt{r2} (line 7).  A sequential proxy will hang up
waiting for the response for \texttt{r1}, never forwarding \texttt{r2}
to the server.  The check on line 9 will therefore fail, because
transaction \texttt{r2} does not complete.  A concurrent proxy can
handle multiple transactions simultaneously, and so it will forward
\texttt{r2} to the server and send the response back to the client, even while waiting for the server to
respond to \texttt{r1}

\begin{figure}
\begin{center}
\input{examples/C01}
\end{center}
\mycaption{Using \pxydrive{} to perform transaction requiring caching.}{A caching proxy will respond to \texttt{r1} without any action by the server.}
\label{cmd:C01}
\end{figure}

Figure \ref{cmd:C01} illustrates the ability of \pxydrive{} to
ascertain that the proxy caches previous results.  It first performs
transaction \texttt{f1}, retrieving file \texttt{random-text.txt}.  It
then issues a request for the same file (transaction \texttt{r1}).  If
the proxy has the file it its cache, it will shortcircuit the
transaction and return the file directly to the client, skipping steps
2--4 of Figure \ref{fig:transaction}.  The server never receives the
request and therefore need not be instructed to respond to it.


\section{Standard Tests}
\label{sec:standard}

The directory \texttt{tests} contains a set of 32 command files
testing many aspects of a proxy.  These are divided into three {\em
  series}, named ``A'',
``B'', and ``C'' based on the first letter of their file names.

The A series (12 files)
tests attributes
required of all proxies, including ones that operate sequentially.  These include the
ability to handle both text and binary files, to correctly handle
exceptional cases, and to support sequences of transactions involving
multiple servers.

The B series (10 files)
tests attributes
required of
a concurrent proxy.  These include cases where the reponses occur in a
different order than the requests, and ones with many transactions,
having the potential to cause race conditions.

The C series (10 files)
tests attributes
required for a caching proxy.  These include the ability to cache text and binary
files, as well as responses for nonexistent files.  It tests the
cache policies on maximum file size, maximum total cache size, and
eviction rules.  These tests require a concurrent
proxy.

Of course, even passing these 32 tests does not ensure complete correctness
of a proxy program.  Errors involving incorrect synchronization are
especially difficult to detect in any testing regime.

\subsection*{Regression Testing with \pxyregress}

The program \pxyregress{} automates the process of running a subset of
the standard tests using \pxydrive{}.
From its local directory, it is invoked as:
\texttt{./pxyregress.py -p} {\it PROXY}, where {\it PROXY} is the path
to the file
containing the executable proxy program.  Command-line options include 
\begin{description}
\item[\texttt{-h}]  $\;$ \\
Print usage information
\item[\texttt{-s} {\it SERIES}] $\;$ \\
Specify which series of tests to
  perform.  By default, all three series (A, B, and C) are tested.  This argument
  allows finer control, specified as a string consisting of some
  combination of the characters
  `\texttt{A}',   `\texttt{B}', and   `\texttt{C}'.  For example, the
  command-line option \texttt{-s A} will run only the A-series tests.
\item[\texttt{-t} {\it SECS}] $\;$ \\
  Set a timeout limit of {\it SECS} seconds for each test.  Any test exceeding
  that limit will be marked as having failed.  The default limit is
  60 seconds.  Setting the limit to 0 indicates that the test should
  be allowed to run indefinitely.
\item[\texttt{-l} {\it FILE}]  $\;$ \\
Copy program outputs to {\it FILE}\@.
\item[\texttt{-S}]  $\;$ \\
Enforce strict formatting rules for HTTP request message headers. 
\end{description}

\pxyregress{} stores the outputs generated by each of the tests in the
\texttt{logs} subdirectory.

\section{Additional Notes}
\label{sec:notes}

\begin{itemize}
\item
The results of \pxydrive{} (and therefore \pxyregress{}) are
nondeterministic, due to the variable delays of network communications
and proxy actions.  From one run to another, a test may pass one time
and fail another.  When running one of the standard tests, such an
inconsistency most likely indicates a synchronization error in the
proxy, since the delays included in the regression tests should
provide ample time for activities to take place.  Keep in mind that a
correct proxy should pass all tests every time.  If tests succeed in
some runs but fail in others, it is highly likely that the proxy is
incorrect, and it is only by chance that the bugs are not always
triggered.

\item
\pxydrive{} uses randomly generated port numbers for the proxy and the
servers.  There is a small chance that its chosen numbers will
conflict with ones already in use.   If this happens, try rerunning
the program, so that it will select a new set of ports.


\item
\pxydrive{} is still in early stages
of development.  It therefore has the potential to hang up when an
error occurs, making it impossible to exit with the \texttt{quit}
command or via Control-C\@.  If this occurs, you can kill the program by
(in a separate window), using the \texttt{ps} command to determine the
process number and the \texttt{kill} command to kill the process.

\end{itemize}

\end{document}
